\subsubsection*{\texttt{52-transform-decomposition}}
\label{app:samples:52_transform_decomposition}

\paragraph{Q3: How long does the Haste spell last?}
\noindent \textbf{Corr.:}~5.00/5\quad \textbf{Faith.:}~1.00\quad \textbf{Relev.:}~1.00\quad \textbf{Ctx.Rel.:}~1.00\\

\noindent\textbf{Antwort:} The Haste spell lasts for a duration of concentration, up to 1 minute.

\noindent\textbf{Abgerufene Chunks:}
\begin{itemize}[noitemsep, topsep=2pt]
  \item[\textbf{\#1}] \texttt{[—]} \textit{?/?} -- Sub question: What is the duration of the Haste spell Response:  The duration of\ldots
  \item[\textbf{\#2}] \texttt{[0.0167]} \textit{Spells/Haste} -- \#\#\# Haste  *3rd-level transmutation*  **Casting Time:** 1 action  **Range:** 30 \ldots
  \item[\textbf{\#3}] \texttt{[0.0164]} \textit{Spells/# Spellcasting} -- \#\# Duration  A spell's duration is the length of time the spell persists. A dura\ldots
  \item[\textbf{\#4}] \texttt{[0.0161]} \textit{Classes/Sorcerer} -- \#\#\#\# Extended Spell  When you cast a spell that has a duration of 1 minute or lo\ldots
  \item[\textbf{\#5}] \texttt{[0.0159]} \textit{Spells/Modify Memory} -- \#\#\# Modify Memory  *5th-level enchantment*  **Casting Time:** 1 action  **Range:\ldots
\end{itemize}
\vspace{4pt}
\hrule
\vspace{6pt}

\paragraph{Q5: Can a grappled creature cast spells?}
\noindent \textbf{Corr.:}~1.00/5\quad \textbf{Faith.:}~0.00\quad \textbf{Relev.:}~0.00\quad \textbf{Ctx.Rel.:}~—\\

\noindent\textbf{Antwort:} The provided context does not specify whether a grappled creature can or cannot cast spells. Therefore, based on the information given, it is unclear if there are any restrictions on a grappled creature's ability to cast spells.

\noindent\textbf{Abgerufene Chunks:}
\begin{itemize}[noitemsep, topsep=2pt]
  \item[\textbf{\#1}] \texttt{[—]} \textit{?/?} -- Sub question: What are the rules regarding grappling in DnD 5e? Response:  To ma\ldots
  \item[\textbf{\#2}] \texttt{[—]} \textit{?/?} -- Sub question: What are the rules for casting spells in DnD 5e? Response:  Accord\ldots
  \item[\textbf{\#3}] \texttt{[—]} \textit{?/?} -- Sub question: Are there any conditions that affect a grappled creature's ability\ldots
  \item[\textbf{\#4}] \texttt{[0.0167]} \textit{Gameplay/Combat} -- \#\#\# Grappling  When you want to grab a creature or wrestle with it, you can use \ldots
  \item[\textbf{\#5}] \texttt{[0.0164]} \textit{Characterizations/Feats} -- \#\# Grappler  *Prerequisite: Strength 13 or higher*  You've developed the skills \ldots
\end{itemize}
\vspace{4pt}
\hrule
\vspace{6pt}

\paragraph{Q12: If a creature has exhaustion level 3 and is also restrained, what is its movement speed?}
\noindent \textbf{Corr.:}~5.00/5\quad \textbf{Faith.:}~1.00\quad \textbf{Relev.:}~1.00\quad \textbf{Ctx.Rel.:}~—\\

\noindent\textbf{Antwort:} The creature's movement speed would be 0 since it is restrained. The exhaustion level does not affect the movement speed directly, as per the provided rules.

\noindent\textbf{Abgerufene Chunks:}
\begin{itemize}[noitemsep, topsep=2pt]
  \item[\textbf{\#1}] \texttt{[—]} \textit{?/?} -- Sub question: What are the effects of exhaustion level 3 on movement speed? Resp\ldots
  \item[\textbf{\#2}] \texttt{[—]} \textit{?/?} -- Sub question: What are the effects of being restrained on movement speed? Respon\ldots
  \item[\textbf{\#3}] \texttt{[0.0167]} \textit{Gamemastering/Conditions} -- \#\# Exhaustion  Some special abilities and environmental hazards, such as starvat\ldots
  \item[\textbf{\#4}] \texttt{[0.0164]} \textit{Gamemastering/Madness} -- **Table- Madness Long-Term Effects**  | d100   | Effect (lasts 1d10 × 10 hours)
  \item[\textbf{\#5}] \texttt{[0.0161]} \textit{Gameplay/Adventuring} -- \#\#\# Travel Pace  While traveling, a group of adventurers can move at a normal, f\ldots
\end{itemize}
\vspace{4pt}
\hrule
\vspace{6pt}

\paragraph{Q16: Can you use Action Surge to cast two leveled spells in one turn?}
\noindent \textbf{Corr.:}~1.00/5\quad \textbf{Faith.:}~0.00\quad \textbf{Relev.:}~0.00\quad \textbf{Ctx.Rel.:}~—\\

\noindent\textbf{Antwort:} Based on the context provided, no, a character cannot use Action Surge to cast two leveled spells in one turn. The Action Surge feature allows a Fighter to take one additional action on top of their regular action, but it does not specify that this additional action can be used to cast two leveled s\ldots

\noindent\textbf{Abgerufene Chunks:}
\begin{itemize}[noitemsep, topsep=2pt]
  \item[\textbf{\#1}] \texttt{[—]} \textit{?/?} -- Sub question: What is Action Surge and how does it work? Response:  Action Surge\ldots
  \item[\textbf{\#2}] \texttt{[—]} \textit{?/?} -- Sub question: What are the rules for casting spells in one turn? Response:  A ch\ldots
  \item[\textbf{\#3}] \texttt{[—]} \textit{?/?} -- Sub question: Are there any limitations on using Action Surge for spellcasting? \ldots
  \item[\textbf{\#4}] \texttt{[—]} \textit{?/?} -- Sub question: What are the requirements for casting leveled spells? Response:  T\ldots
  \item[\textbf{\#5}] \texttt{[0.0167]} \textit{Classes/Fighter} -- \#\#\# Action Surge  Starting at 2nd level, you can push yourself beyond your norma\ldots
\end{itemize}
\vspace{4pt}
\hrule
\vspace{6pt}
